% Part: proof-theory
% Chapter: sequent-calculus
% Section: interpretation-rules

\documentclass[../../../include/open-logic-section]{subfiles}

\begin{document}

\olfileid{pt}{seq}{irl}

\olsection{تفسیر قواعد}

تفسیر یک دنباله را به یاد آورید: دنبالهٔ $\Gamma \Sequent
\Delta$ در !!a{structure}~$\Struct{M}$ صادق است هرگاه دست‌کم یک
$!A \in \Gamma$ کاذب باشد یا دست‌کم یک~$!A \in \Delta$ صادق باشد. می‌توان
قواعد منطقیِ~\Log{G3c} را صورت‌بندیِ شرایط صدقِ !!{formula}‌های اصلیِ
آن‌ها دانست. قاعدهٔ \RightR{\lif} را در نظر بگیرید. !!{formula}ٔ اصلیِ
آن $!A \lif !B$ است. این فرمول هنگامی صادق است که یا $!A$ کاذب باشد یا
~$!B$ صادق باشد. ازاین‌رو، دنبالهٔ $\ \Sequent !A \lif !B$ صادق است اگر
و تنها اگر $!A \Sequent !B$ صادق باشد. اگر !!{formula}‌های بافتِ
$\Gamma$ و~$\Delta$ را به‌ترتیب به مقدم و تالیِ این دنباله‌ها بیفزاییم،
دقیقاً نتیجه و مقدمهٔ قاعدهٔ \RightR{\lif} به دست می‌آید. وضع برای
~\LeftR{\lif} نیز مشابه است. $!A \lif !B$ هنگامی کاذب است که $!A$ صادق و
$!B$ کاذب باشد. بنابراین، $!A \lif !B \Sequent \ $ صادق است اگر و تنها
اگر هر دو دنبالهٔ $\ \Sequent !A$ و $!B \Sequent \ $ صادق باشند. نتیجهٔ
\LeftR{\lif} هم‌ارزِ عطفِ مقدمات آن است. این الگو تضمین می‌کند که قواعد
~\Log{G3c} درست‌اند (اگر همهٔ مقدمات صادق باشند، نتیجه نیز صادق است).
همچنین تمامیت را تضمین می‌کند.

در واقع، شرایط صدقِ هر رابطِ تابع‌ارزشی را می‌توان بدین شیوه با مجموعه‌ای
از دنباله‌ها توصیف کرد. دلیل آن این است که هر تابع صدق در منطق کلاسیک را
می‌توان با !!a{formula} در صورت نرمال عطفی توصیف کرد: عطفی از فصل‌های
!!{formula}‌های اتمی و نقیضِ اتمی. برای نمونه، فرض کنید $!A \oplus !B$
هنگامی صادق باشد که دقیقاً یکی از $!A$ و $!B$ صادق باشد؛ یعنی «یای
انحصاری». آنگاه $!A \oplus !B$ صادق است اگر و تنها اگر $(!A \lor !B)
\land (\lnot !A \lor \lnot !B)$ صادق باشد، و هنگامی کاذب است که
$(\lnot !A \lor !B) \land (!A \lor \lnot !B)$ صادق باشد. پس می‌توان
قواعد حساب دنباله‌ایِ $\oplus$ را چنین معرفی کرد:
\begin{align*}
&
\Axiom$\Gamma \fCenter \Delta, !A, !B$
\Axiom$!A, !B, \Gamma \fCenter \Delta$
\RightLabel{\RightR{\oplus}}
\BinaryInf$\Gamma \fCenter \Delta, !A \oplus !B$
\DisplayProof
\\
& \Axiom$!A, \Gamma \fCenter \Delta, !B$
\Axiom$!B, \Gamma \fCenter \Delta, !A$
\RightLabel{\LeftR{\oplus}}
\BinaryInf$!A \oplus !B, \Gamma \fCenter \Delta$
\DisplayProof
\end{align*}
این قواعد نیز درست و کامل خواهند بود. نتیجهٔ هر قاعده صادق است اگر و تنها
اگر همهٔ مقدمات آن صادق باشند.

\begin{prob}
فرض کنید $!A \downarrow !B$ هنگامی صادق باشد که نه $!A$ و نه~$!B$ صادق
نباشند («NOR» یا «نقیضِ یا»). قواعد آن در \Log{G3c} چه خواهند بود؟
(دو دنباله بیابید که هم‌زمان صادق‌اند اگر و تنها اگر $!A \downarrow !B$
صادق باشد؛ این دو، قاعدهٔ \RightR{\downarrow} را به دست می‌دهند. سپس یک
دنباله بیابید که صادق است اگر و تنها اگر $!A \downarrow !B$ کاذب باشد؛
این دنباله قاعدهٔ \LeftR{\downarrow} را به دست می‌دهد.)
\end{prob}

در \Log{G3c}، هر رابط و هر سور دقیقاً دو قاعده دارد: یک قاعدهٔ چپ و یک
قاعدهٔ راست. اما در \Log{G1c} دو قاعدهٔ \LeftR{\land} و دو قاعدهٔ
\RightR{\lor} وجود دارد. علت آن است که دستگاه اصلی گنتسن، یعنی~\Log{LK}
که \Log{G1c} بر الگوی آن ساخته شده است، گونه‌ای به نام~\Log{LJ} برای
منطق غیرکلاسیک مهمی، یعنی منطق شهودی، داشت. برای به‌دست‌آوردن دستگاهی
درست و کامل برای منطق شهودی، تالیِ هر دنباله در~\Log{LJ} به داشتن حداکثر
یک !!{formula} محدود می‌شود. ازاین‌رو نمی‌توان قاعدهٔ \RightR{\lor} در
~\Log{G3c} را به کار برد، زیرا تالیِ مقدمهٔ آن هم $!A$ و هم~$!B$ را در بر
دارد. بنابراین گنتسن برای \Log{LJ} یک جفت قاعدهٔ \RightR{\lor} به کار
برد و همان‌ها را برای~\Log{LK} نیز برگزید. نسخهٔ شهودی~\Log{G1i} نیز
همان~\Log{G1c} است با این محدودیت که همهٔ تالی‌ها حداکثر یک !!{formula}
داشته باشند. (\Log{G3c} نیز نسخه‌ای شهودی دارد، اما برخی قواعد آن با
قواعد متناظر در~\Log{G3c} متفاوت‌اند؛ برای مثال، آن دستگاه نیز از یک
جفت قاعدهٔ \RightR{\lor} استفاده می‌کند.)

درستیِ قواعد ساختاریِ~\Log{G1c} به‌آسانی دیده می‌شود. برای مثال، اگر
$\Gamma \Sequent \Delta$ !!{formula}ی کاذب در سمت چپ یا !!{formula}ی
صادق در سمت راست داشته باشد، $\Gamma \Sequent \Delta, !A$ نیز همان
فرمول را دارد؛
صادق یا کاذب‌بودن $!A$ در این استدلال اثری ندارد. پس اگر مقدمهٔ
~\RightR{\Weakening} صادق باشد، نتیجهٔ آن نیز صادق است. توجه کنید که عکس
این گزاره درست نیست. ممکن است نتیجه به سبب صادق‌بودن $!A$ صادق باشد؛ در
این صورت صادق‌بودن مقدمه تضمین نمی‌شود.

اصلاً چرا قواعد ساختاری داشته باشیم؟ برای نمونه، انقباض
(\RightR{\Contraction} و \LeftR{\Contraction}) لازم است تا بتوان
$\ \Sequent !A \lor \lnot !A$ را !!{prove} کرد. اگر از $!A \Sequent !A$
آغاز کنید، نخست با~\RightR{\lnot} به $\Sequent !A, \lnot !A$ می‌رسید.
سپس می‌توانید دو نسخهٔ \RightR{\lor} را پی‌درپی به کار ببرید: یک بار
برای به‌دست‌آوردن $!A \lor \lnot !A$ از $\lnot !A$ و بار دیگر برای
به‌دست‌آوردن $!A \lor \lnot !A$ از~$!A$. حاصل
$\ \Sequent !A \lor \lnot !A, !A \lor \lnot !A$ است. بنابراین، برای
رسیدن به $\Sequent !A \lor \lnot !A$ در~\Log{G1c} باید بتوان دو رخداد
از یک !!{formula} را در !!a{proof} «منقبض» کرد:
\[
\Axiom$!A \fCenter !A$
\RightLabel{\RightR{\lnot}}
\UnaryInf$\fCenter !A, \lnot !A$
\RightLabel{\RightR{\lor}}
\UnaryInf$\fCenter !A, !A \lor \lnot !A$
\RightLabel{\RightR{\lor}}
\UnaryInf$\fCenter !A \lor \lnot !A, !A \lor \lnot !A$
\RightLabel{\RightR{\Contraction}}
\UnaryInf$\fCenter !A \lor \lnot !A$
\DisplayProof
\]
در واقع، دنباله‌های معتبری وجود دارند که بدون تضعیف یا انقباض نمی‌توان
آن‌ها را در~\Log{G1c} !!{prove} کرد. برای مثال، !!a{proof} از دنبالهٔ
$!A \Sequent !B \lif !A$ در~\Log{G1c} باید به \RightR{\lif} ختم شود
(زیرا تنها $!B \lif !A$ می‌تواند اصلی باشد)، و این امر دنبالهٔ
$!B, !A \Sequent !A$ را به‌عنوان مقدمه می‌طلبد. این مقدمه نیز به
\LeftR{\Weakening} نیاز دارد تا از اصل $!A \Sequent !A$ به دست آید:
\[
\Axiom$!A \fCenter !A$
\RightLabel{\LeftR{\Weakening}}
\UnaryInf$!B, !A \fCenter !A$
\RightLabel{\RightR{\lif}}
\UnaryInf$!A \fCenter !B \lif !A$
\DisplayProof
\]

در \Log{G3c} اصلاً به قواعد انقباض و تضعیف نیاز نیست. تضعیف در خود اصول
گنجانده شده است. از انقباض نیز عمدتاً می‌توان با یکی‌کردن دو نسخهٔ بافتِ
$\Gamma$ و~$\Delta$ در قواعد دومقدمه‌ای پرهیز کرد. هم \Log{G1c} و هم
\Log{G3c} چنین قواعد «اشتراک بافتی» دارند. در قواعد تک‌مقدمه‌ای، انقباض
فقط برای \RightR{\lor}، \LeftR{\land}، \RightR{\lexists} و
~\LeftR{\lforall} لازم است. صورت این قواعد در~\Log{G3c} انقباض را کاملاً
زائد می‌کند. دستگاهی به نام \Log{G2c} نیز وجود دارد (نگاه کنید به
\olref{tab:G2c}) که در آن قواعد دومقدمه‌ای بافت‌های مستقلی دارند. در
\Log{G2c} انقباض حتی از~\Log{G1c} نیز مهم‌تر است.

\subfile{rules-G2c}

\end{document}
